\NativeSlot{43}{87}{%
% p090_s01

\ExamQuestion{3.}{5}%
气象观测站在某天的\(3\)时、\(8\)时、\(15\)时、\(21\)时测得的温度分别是\(8^\circ\mathrm C\)、\(15^\circ\mathrm C\)、\(26^\circ\mathrm C\)、\(19^\circ\mathrm C\)。计算这天的平均气温正确的算式为\mbox{（\hspace{1.8em}）}。\par
\ExamOptions{A. \((8+15+26+19)\div4\)}{B. \((8+15+26+19)\div(3+8+15+21)\)}{C. \((8+15+26+19)\div24\)}{D. \((3+8+15+21)\div4\)}
}
\NativeSlot{194}{86}{%
% p090_s02

\ExamQuestion{4.}{5}%
一次跳绳比赛，\allowbreak{}\(5\)名同学的成绩从低到高依次是\(73\)下、\(81\)下、\(a\)下、\(87\)下、\(95\)下，他们的平均成绩四舍五入到个位，可能是\mbox{（\hspace{1.8em}）}。\par
\ExamOptions{A. 72下}{B. 85下}{C. 95下}{D. 96下}
}
\NativeSlot{344}{87}{%
% p090_s03

\ExamQuestion{5.}{5}%
\(5\)个数的平均数是\(50\)，若把其中一个数改为\(80\)，则平均数变为\(60\)，这个数原来是\mbox{（\hspace{1.8em}）}。\par
\ExamOptions{A. 60}{B. 50}{C. 40}{D. 30}
}
\NativeSlot{497}{56}{%
\ExamGroup{三、计算。{\normalfont\kaishu\fontsize{10.8bp}{14bp}\selectfont （共23分）}}

% p090_s04

\ExamQuestion{1.}{5}%
直接写得数。（8分）\par
\noindent \NativeOralCell{0.4+0.5=}\NativeOralCell{1-0.7=}\NativeOralCell{4.3+3.7=}\NativeOralCell{1.2-0.8=}\par

% p090_s05

\noindent \NativeOralCell{0.4+1.2=}\NativeOralCell{7.9-1.9=}\NativeOralCell{7.3+1.8=}\NativeOralCell{13.5-3.5=}\par
}
\NativeSlot{596}{117}{%
\noindent \textbf{2.}\quad 列竖式计算。（带☆的要验算）（15分）\par\smallskip
\noindent%
% p090_s06
\makebox[.249\linewidth][l]{\(15.05+2.63=\hspace{1.8em}\)}%
% p090_s07
\makebox[.249\linewidth][l]{\(9.06-2.5=\hspace{1.8em}\)}%
% p090_s08
\makebox[.249\linewidth][l]{\(5.7+12.34=\hspace{1.8em}\)}%
% p090_s09
\makebox[.249\linewidth][l]{☆\(12.8-1.24=\hspace{1.8em}\)}\par
}
